\begin{figure}[hbt]
\centering
\begin{subfigure}{0.3\textwidth}
\centering
\begin{tikzpicture}[scale=0.7]
    \draw[black, thick, fill=ApplePalette1!50] (0,0,4) -- (2,0,4) -- (2,4,4) -- (0,4,4) -- cycle;
    \draw[black, thick, fill=ApplePalette4!50] (2,0,4) -- (4,0,4) -- (4,4,4) -- (2,4,4) -- cycle;% front
    \draw[black, thick, fill=ApplePalette1!50] (0,4,0) -- (2,4,0) -- (2,4,4) -- (0,4,4) -- cycle; % top
    \draw[black, thick, fill=ApplePalette4!50] (2,4,0) -- (4,4,0) -- (4,4,4) -- (2,4,4) -- cycle; % top
    \draw[black, thick, fill=ApplePalette4!50] (4,0,0) -- (4,4,0) -- (4,4,4) -- (4,0,4) -- cycle; % right
\end{tikzpicture}
\caption{Slab decomposition.}
\end{subfigure}
\hfill
\begin{subfigure}{0.3\textwidth}
\centering
\begin{tikzpicture}[scale=0.7]
    \draw[black, thick, fill=ApplePalette1!50] (0,0,4) -- (2,0,4) -- (2,2,4) -- (0,2,4) -- cycle; % front
    \draw[black, thick, fill=ApplePalette2!50] (2,0,4) -- (4,0,4) -- (4,2,4) -- (2,2,4) -- cycle; % front
    \draw[black, thick, fill=ApplePalette4!50] (0,2,4) -- (2,2,4) -- (2,4,4) -- (0,4,4) -- cycle; % front
    \draw[black, thick, fill=ApplePalette6!50] (2,2,4) -- (4,2,4) -- (4,4,4) -- (2,4,4) -- cycle; % front
    \draw[black, thick, fill=ApplePalette4!50] (0,4,0) -- (2,4,0) -- (2,4,4) -- (0,4,4) -- cycle; % top
    \draw[black, thick, fill=ApplePalette6!50] (2,4,0) -- (4,4,0) -- (4,4,4) -- (2,4,4) -- cycle; % top
    \draw[black, thick, fill=ApplePalette6!50] (4,2,0) -- (4,4,0) -- (4,4,4) -- (4,2,4) -- cycle; % right
    \draw[black, thick, fill=ApplePalette2!50] (4,0,0) -- (4,2,0) -- (4,2,4) -- (4,0,4) -- cycle; % right
\end{tikzpicture}
\caption{Pencil decomposition.}
\end{subfigure}
\hfill
\begin{subfigure}{0.3\textwidth}
\centering
\begin{tikzpicture}[scale=0.7]
    \draw[black, thick, fill=ApplePalette1!50] (0,0,4) -- (2,0,4) -- (2,2,4) -- (0,2,4) -- cycle; % front
    \draw[black, thick, fill=ApplePalette2!50] (2,0,4) -- (4,0,4) -- (4,2,4) -- (2,2,4) -- cycle; % front
    \draw[black, thick, fill=ApplePalette4!50] (0,2,4) -- (2,2,4) -- (2,4,4) -- (0,4,4) -- cycle; % front
    \draw[black, thick, fill=ApplePalette6!50] (2,2,4) -- (4,2,4) -- (4,4,4) -- (2,4,4) -- cycle; % front
    \draw[black, thick, fill=ApplePalette4!50] (0,4,2) -- (2,4,2) -- (2,4,4) -- (0,4,4) -- cycle; % top
    \draw[black, thick, fill=ApplePalette6!50] (2,4,2) -- (4,4,2) -- (4,4,4) -- (2,4,4) -- cycle; % top
    \draw[black, thick, fill=ApplePalette6!50] (4,2,2) -- (4,4,2) -- (4,4,4) -- (4,2,4) -- cycle; % right
    \draw[black, thick, fill=ApplePalette2!50] (4,0,2) -- (4,2,2) -- (4,2,4) -- (4,0,4) -- cycle; % right
    \draw[black, thick, fill=ApplePalette5!50] (0,4,0) -- (2,4,0) -- (2,4,2) -- (0,4,2) -- cycle; % top
    \draw[black, thick, fill=ApplePalette3!50] (2,4,0) -- (4,4,0) -- (4,4,2) -- (2,4,2) -- cycle; % top
    \draw[black, thick, fill=ApplePalette3!50] (4,2,0) -- (4,4,0) -- (4,4,2) -- (4,2,2) -- cycle; % right
    \draw[black, thick, fill=ApplePalette7!50] (4,0,0) -- (4,2,0) -- (4,2,2) -- (4,0,2) -- cycle; % right
    \end{tikzpicture}
\caption{Brick decomposition.}
\end{subfigure}
\caption{Domain decomposition strategies for parallel computation. Separate colours refer to separate device partitions. Tikz source code: Tikz/decompStrategies.tex}\label{fig:decompStrategies}
\end{figure}%